\documentclass[10pt,a4paper]{amsart}
\usepackage[margin=19mm]{geometry}
\usepackage{amsmath,amssymb,mathtools}
\usepackage[scr=rsfs,cal=euler]{mathalfa}
\usepackage{xfrac}
\usepackage[unicode,hidelinks,bookmarks=false]{hyperref}
\newcommand{\ie}{i.e.\ }
\newcommand{\cf}{cf.\ }
\newcommand{\resp}{respectively\ }
\newcommand{\Subsection}{\subsection}
\providecommand{\nobreakdash}{\nobreak\hbox{-}}
\setlength{\emergencystretch}{2em}
\multlinegap0pt
\usepackage{bm}
\usepackage{bbm}
\usepackage{dsfont}
\usepackage{xspace}
\usepackage{cancel}
\usepackage{mathtools}
\usepackage{cleveref}
\usepackage[shortlabels]{enumitem}
\usepackage{amsthm}
\makeatletter
\@ifundefined{thm}{\newtheorem{thm}{Thm}}{}
\@ifundefined{claim}{\newtheorem{claim}[thm]{Claim}}{}
\@ifundefined{lem}{\newtheorem{lem}[thm]{Lemma}}{}
\@ifundefined{proposition}{\newtheorem{proposition}[thm]{Proposition}}{}
\@ifundefined{rem}{\newtheorem{rem}[thm]{Remark}}{}
\makeatother
\newcommand{\ram}{{\rm ram}\,}
\title[Hyperbolicity and fundamental groups of complex quasi-projec]{Hyperbolicity and fundamental groups of complex quasi-projective varieties (I): Maximal quasi-Albanese dimension by Nevanlinna theory}
\author{Benoit Cadorel, Ya Deng, Katsutoshi Yamanoi}
\date{Source: arXiv:2511.04405v1 (CC BY 4.0)}
\begin{document}
\maketitle

\noindent\emph{Source and licence.} Hyperbolicity and fundamental groups of complex quasi-projective varieties (I): Maximal quasi-Albanese dimension by Nevanlinna theory. Version arXiv:2511.04405v1 (CC BY 4.0), distributed under the Creative Commons Attribution 4.0 International licence, as recorded in the arXiv licence metadata for this exact version and shown on the abstract page https://arxiv.org/abs/2511.04405v1. Full rights notice: licenses/arxiv-2511.04405-CC-BY-4.0.txt.\par\smallskip

\noindent\textit{Editorial scope.} Counted unit: the Proposition whose source label is \texttt{prop:20230405} in the article (source0.tex lines 3444--3457, source label \texttt{prop:20230405}), delivered together with the article's own complete original proof of it (source0.tex lines 3459--3561, one \texttt{proof} environment of 103 lines). No mathematics is written, repaired or re-derived: statement and proof are the article's own text line for line. Required context retained uncounted: lem:20230411 (lines 980--984); rem:20221203 (lines 2683--2688); lem:20220909gg (lines 2698--2719). Every definition, display and label of those lines is retained unchanged. Omitted from this package and not redistributed: 1--979, 985--2682, 2689--2697, 2720--3443, 3458--3458, 3562--4059. Nothing inside a retained excerpt is omitted or altered: every retained body is delivered exactly as the article's lines are written, so no omission receipt of any kind accompanies this package. Citations: the retained excerpts carry no citation command at all, so no bibliography entry of the article is transcribed and no reference is redistributed. Reference adaptation: none was needed and none was made; every reference inside the delivered unit resolves inside it, so no unresolved reference or citation is produced. Printed slips kept literal: no printed word, symbol, hypothesis or number of the counted statement or of its proof was altered. Layout and class: the article's own class is \texttt{smfart} with packages including bm, babel, amscd, amssymb, url, xspace, smfthm, paralist, datetime, colortbl, color, xcolor, hyperref, mathtools; the wrapper is the lane's pinned \texttt{amsart} editorial wrapper at 10pt on A4, which restates the theorem-like environments the retained text needs and adds the article's own macro definitions that the retained text needs (\texttt{\textbackslash ram}), with extra wrapper packages (bm, bbm, dsfont, xspace, cancel, mathtools, cleveref, enumitem). The delivered numbering is the unit's own. Assets: no retained span contains a figure environment, a graphics-inclusion command, a PDF-page inclusion, a table or a TikZ picture; no image file is copied into this package, the package declares no asset dependency and no image file is redistributed with it. Licence: version arXiv:2511.04405v1 of this article is distributed under the Creative Commons Attribution 4.0 International licence, as recorded in the arXiv licence metadata for this exact version and shown on the abstract page https://arxiv.org/abs/2511.04405v1. A full rights notice is delivered at licenses/arxiv-2511.04405-CC-BY-4.0.txt. Characters: every retained excerpt is pure ASCII, so the delivered engine, XeLaTeX with its default Latin Modern text font, reports no missing character.
	\begin{lem}\label{lem:20230411}
		Let $X$ be a smooth projective variety and let $f:Y\to X$ be a holomorphic map.
		Assume that $T_f(r)=O(\log r)||$ and $N_{\ram\pi}(r)=O(\log r)||$. 
		Then $f$ does not have essential singularity over $\infty$.
	\end{lem}

	\begin{rem}\label{rem:20221203}
		If $f\in I_{\{0\}}$, then $f:Y\to A\times S$ does not have essential singularity over $\infty$.
		Indeed, $f\in I_{\{0\}}$ yields $T_{f_A}(r)=O(\log r)+o(T_f(r))||$ and  $T_{f_S}(r)=O(\log r)+o(T_f(r))||$.
		Hence $T_f(r)=O(\log r)||$, thus $N_{\ram\pi}(r)=O(\log r)||$.
		These two estimates implies that $f$ does not have essential singularity over $\infty$ (cf. \cref{lem:20230411} ). 
	\end{rem}

	\begin{lem}\label{lem:20220909gg}
Let $B \subset A$ be a semi-abelian subvariety, and let $W \subset A \times S$ be a $B$-invariant closed subvariety, where $S$ is projective. Let $D \subset W$ be a reduced Weil divisor.  
Then the following objects exist with the specified property described later: 
\begin{itemize}
    \item a smooth projective equivariant compactification $A \subset \overline{A}$,
    \item a proper birational morphism $\varphi: V \to \overline{W}$ with $V$ smooth, such that 
    \(\varphi^{-1}(\overline{D} \cup \partial W) \subset V\) is a simple normal crossing divisor.
\end{itemize}
The satisfied property: Let $Z \subset W$ be a Zariski closed subset with $\mathrm{codim}(Z,W) \ge 2$ and $Z \subset D$, let $L$ be an ample line bundle on $V$, and let $\varepsilon > 0$. Then there exist
\begin{itemize}
    \item a finite subset $P \subset \mathcal{S}(B) \setminus \{B\}$, and
    \item a proper Zariski closed subset $\Phi \subsetneqq W$ with $\varphi(\mathrm{Ex}(\varphi)) \subset \overline{\Phi}$,
\end{itemize}
such that for every holomorphic map $f: Y \to W$ satisfying $f(Y) \not\subset D \cup \Phi$ and $f \in I_B$, one of the following holds:
\begin{enumerate}
    \item There exists $C \in P$ such that $f \in I_C$.
    \item Let $f': Y \to V$ be the lift of $f$. Then
    \[
    T_{f'}(r, K_V(\varphi^{-1}(\overline{D} \cup \partial W))) \le \overline{N}_f(r, D \setminus Z) + \varepsilon\, T_{f'}(r,L) + O(\log r) + o(T_{f'}(r))||.
    \]
\end{enumerate}
\end{lem}

	\begin{proposition}\label{prop:20230405}
		Let $D\subset A$ be a reduced divisor such that $\mathrm{St}(D)=\{ a\in A; a+D=D\}$ is finite.
		Let $Z\subset D$ be a Zariski closed subset such that $\mathrm{codim}(Z,A)\geq 2$.
		Then there exists a proper Zariski closed subset $\Xi\subsetneqq A$ with the following property:
		Let $f:Y\to A$ be a holomorphic map with the following three properties:
		\begin{enumerate}[label=(\alph*)]
			\item  
			$N_{\ram\pi}(r)=O(\log r)+o(T_f(r))||$, 
			\item  
			$f(Y)\not\subset \Xi\cup D$,
			\item
			$\mathrm{ord}_yf^*D\geq 2$ for all $y\in f^{-1}(D\backslash Z)$.	\end{enumerate}
		Then $f$ does not have essential singularity over $\infty$.
	\end{proposition}

	\begin{proof}
		For a semi-abelian variety $B\subset A$, we denote by $J_B$ the set of all holomorphic maps $f:Y\to A$ such that 
		\begin{itemize}
			\item
			$f\in I_B$,  
			\item
			$f(Y)\not\subset D$, and
			\item
			$\mathrm{ord}_yf^*D\geq 2$ for all $y\in f^{-1}(D\backslash Z)$. \end{itemize}
		Given a semi-abelian variety $B\subset A$, we first prove the following claim.
		
		\begin{claim}\label{claim:20230405}
			There exist a finite subset $P_B\subset \mathcal{S}(B)\backslash\{B\}$ and a proper Zariski closed set $\Phi_B\subsetneqq A$ with the following property:
			Let $f:Y\to A$ be a holomorphic map such that $f\in J_B$ and $f(Y)\not\subset \Xi_B$. 
			Then either one of the followings holds:
			\begin{enumerate}[label=(\alph*)]
				\item
				There exists $C\in P_B$ such that $f\in J_C$.
				\item
				$f$ does not have essential singularity over $\infty$.
			\end{enumerate}
		\end{claim}
		
		\begin{proof}[Proof of \cref{claim:20230405}]
			Since $\mathrm{St}(D)$ is finite, $A\backslash D$ is of log-general type.
			We apply \cref{lem:20220909gg} to get a smooth projective equivariant compactification $\overline{A}$ and a proper birational morphism $\varphi:V\to \overline{A}$, where $V$ is smooth and $\varphi^{-1}(\overline{D}\cup \partial A)$ is a simple normal crossing divisor.
			Since $A\backslash D$ is of log-general type, we deduce that $\varphi ^{-1}(A\backslash D)=V\backslash \varphi^{-1}(\overline{D}\cup \partial A)$ is also of log-general type.
			Thus $K_{V}(\varphi^{-1}(\overline{D}\cup \partial A))$ is big.
			
			By Kodaira's lemma, there exist an effective divisor $E\subsetneqq \overline{A}'$ and a positive integer $l\in\mathbb Z_{\geq 1}$ such that $L=lK_{V}(\varphi^{-1}(\overline{D}\cup \partial A))-E$ is ample.	
			Hence if $g:Y\to V$ satisfies $g(Y)\not\subset E$, then 
			$$T_{g}(r)= O( T_{g}(r,K_{V}(\varphi^{-1}(\overline{D}\cup \partial A))))+O(\log r).$$ 	
			
			
			
			By \cref{lem:20220909gg} applied to $(Z\cup \varphi (\mathrm{Ex}(\varphi)))\cap D$, $L$ and $\varepsilon=1/3l$, we get a finite subset $P_B\subset \mathcal{S}(B)\backslash\{B\}$ and a proper Zariski closed set $\Phi_B\subsetneqq A$ with $\varphi (\mathrm{Ex}(\varphi))\subset \overline{\Phi_B}$.
			We set $\Xi_B=\Phi_B\cup\varphi(E)$.
			Let $f:Y\to A$ be a holomorphic map such that $f\in J_B$ and $f(Y)\not\subset \Xi_B$.
			Then $ f(Y)\not\subset D\cup \Xi_B$ and $f\in I_B$.
			Suppose that the first assertion of \cref{claim:20230405} is not valid.
			Then $f\not\in I_C$ for all $C\in P_B$.
			By \cref{lem:20220909gg}, we get
			\begin{equation*}
				T_{g}(r,K_{V}(\varphi^{-1}(\overline{D}\cup \partial A)))
				\leq \overline{N}_{f}(r,D\backslash (Z\cup \varphi (\mathrm{Ex}(\varphi))))+\frac{1}{3l}T_{g}(r,L)
				\\
				+O(\log r)+o(T_{g}(r))||,
			\end{equation*}
			where $g:Y\to V$ is the lift of $ f:Y\to A$.
			By $g(Y)\not\subset E$, we have 
			$$
			T_{g}(r,L)\leq lT_{g}(r,K_{V}(\varphi^{-1}(\overline{D}\cup \partial A)))+O(\log r).
			$$
			Hence we get
			\begin{equation*}
				\frac{2}{3}T_{g}(r,K_{V}(\varphi^{-1}(\overline{D}\cup \partial A)))
				\leq \overline{N}_{f}(r,D\backslash (Z\cup \varphi (\mathrm{Ex}(\varphi))))
				\\
				+O(\log r)+o(T_{g}(r))||.
			\end{equation*}
			
			
			We estimate the first term of the right hand side.
			Since $\overline{A}$ is smooth and equivariant, $\partial A$ is a simple normal crossing divisor.
			Hence we may decompose as $\varphi^{-1}(\overline{D}\cup \partial A)=H+F$ so that $F=\varphi^{-1}(\partial A)$.
			Then $H=\overline{\varphi^{-1}(D)}$.
			The induced map $V\backslash (F\cup \mathrm{Ex}(\varphi))\to A\backslash \varphi(\mathrm{Ex}(\varphi))$ is isomorphic.
			Hence by $\mathrm{ord}_yf^*D\geq 2$ for all $y\in f^{-1}(D\backslash Z)$, we have
			$$
			2\overline{N}_{f}(r,D\backslash (Z\cup \varphi (\mathrm{Ex}(\varphi))))\leq T_g(r,H).
			$$
			By $\bar{\kappa}(V\backslash F)=0$ and $g(Y)\not\subset \mathrm{Ex}(\varphi)$, we have $T_g(r,K_V(F))+O(\log r)>0$.
			Hence
			$$
			2\overline{N}_{f}(r,D\backslash (Z\cup \varphi (\mathrm{Ex}(\varphi))))\leq T_g(r,K_V(H+F))+O(\log r).
			$$
			Hence we get
			$$
			T_{g}(r,K_{V}(H+F))=O(\log r)+o(T_{g}(r))||.
			$$ 	 	
			Thus $T_{g}(r)=O(\log r) ||$.
			This shows $T_{f}(r)=O(\log r) ||$. 
			Hence $f\in I_{\{0\}}$.
			Hence by \cref{rem:20221203}, $f$ does not have essential singularity over $\infty$.
		\end{proof}
		
		
		
		
		We take a sequence $\mathcal{P}_i\subset \mathcal{S}(A)$ of finite sets as follows.
		Set $\mathcal{P}_1=\{A\}$.
		Given $\mathcal{P}_i$, we define $\mathcal{P}_{i+1}$ as follows.
		For each $B\in\mathcal{P}_i$, we apply \cref{claim:20230405} to get $\Xi_B\subsetneqq A$ and $P_B\subset \mathcal{S}(B)\backslash\{B\}$.
		We set $\mathcal{P}_{i+1}=\cup_{B\in \mathcal{P}_i}P_B$ and $\Xi_{i+1}=\cup_{B\in\mathcal{P}_i}\Xi_B$.
		Set $\Xi=\cup_i\Xi_i$.
		Then $\Xi\subsetneqq A$ is a proper Zariski closed set.
		
		Let $f:Y\to A$ satisfy the three properties in \cref{prop:20230405}.
		Then $f\in J_A$.
		We take $i$ which is maximal among the property that there exists $B\in \mathcal{P}_i$ such that $f\in J_B$. 	
		Then by $f(Y)\not\subset \Xi_B\subset \Xi$,
		\cref{claim:20230405} implies that $f$ does not have essential singularity over $\infty$.
	\end{proof}

\end{document}
